\section{Open questions (5 concrete probes)}

\begin{enumerate}
\item \textbf{Direct-sparse vs.\ LCU-of-Paulis baseline on the SAME instance.}
Does the paper's direct sparse-matrix block-encoding actually beat a generic
Jordan--Wigner\,$\to$\,LCU-of-Paulis baseline on the same 3-nucleon /
6-single-particle $H_{\text{pair}}$, in both sub-normalization $\alpha$ and
post-transpilation Clifford+$T$ gate count?
\emph{Basis:} we verified only that the paper's own construction hits
$\alpha=16$ at machine precision; the structural argument (direct beats
generic) is not independently cross-checked in our run.
\emph{Next steps:} JW-transform $H_{\text{pair}}$ via OpenFermion, decompose to
Pauli strings, build a PREP+SEL LCU block-encoding in Qiskit, transpile both
circuits to Clifford+$T$, and compare
$(\alpha,\text{T-count},\text{2q-count})$ at $L{=}9,12,16,20$.

\item \textbf{Empirical $O(L\log L)$ scaling of $O_C$.}
Does the $\approx 12L\log L + 23L$ two-qubit / $14L\log L + 21L$ $T$-gate
scaling (Sec.~4.4) hold after compilation and optimization?
\emph{Basis:} we evaluated the paper's closed-form formula at $L{=}9$ but never
instantiated $O_C$ as a real \texttt{QuantumCircuit} and transpiled it.
\emph{Next steps:} implement $O_C(l)$ as a Qiskit circuit with multi-controlled
swaps, transpile with T-count optimization (pyzx or
\texttt{optimization\_level=3}), measure 2q + $T$ gates at
$L\in\{6,9,12,16,20,30,50\}$, and fit $a + bL\log L + cL$; check residuals
against the paper's constants.

\item \textbf{QSVT / qubitization query cost for a physical DoS task.}
Applied inside a QSVT pipeline (Sec.~5.3), does $\alpha=16$ leave the query
count for a nuclear-structure DoS or ground-state task in the practical
regime?
\emph{Basis:} C5 was flagged out-of-scope for the one-shot replication.
\emph{Next steps:} use qsppack / pyqsp to build a Gaussian-smoothed-DoS QSVT
polynomial, wrap $U_H$ as the walk operator, measure $Q(\alpha,\varepsilon)$
against the $\alpha=16$ prediction, and repeat on a 5-nucleon / 10-basis
instance small enough for classical ground-truth.

\item \textbf{End-to-end fault-tolerant resource estimate.}
Under surface-code assumptions ($d\approx 15$--$25$, $p\approx 10^{-3}$, magic-
state distillation), does the $O(L\log L)$ advantage over LCU persist in
physical-qubit and wall-clock cost?
\emph{Basis:} the current replication is a purely algebraic classical
verification --- no noise model, no FT overhead accounting.
\emph{Next steps:} feed the compiled $O_C$ circuit into the Azure Quantum
Resource Estimator (or \texttt{qsharp-community/resource-estimator}) for a
QSVT-DoS run at $\sigma \sim 0.1$~MeV, comparing to the LCU baseline of Q1.

\item \textbf{ML-selected block-encoding scheme.}
Can an RL policy over Camps--Grover--Rossi--Van~Beeumen decomposition choices
(or a differentiable circuit search) discover a block-encoding of
$H_{\text{pair}}$ with strictly smaller $\alpha$ or gate count than the hand-
designed construction, and does it generalize to related sparse fermionic
Hamiltonians (Richardson, Lipkin--Meshkov--Glick, generalized-seniority
pairing)?
\emph{Basis:} the paper's construction is human-crafted; optimality is not
addressed.
\emph{Next steps:} build an RL environment whose actions are Camps-et-al.\
decomposition primitives (controlled-swap patterns, diffusion register sizes,
ancilla allocations) and whose reward is
$-(\alpha + \lambda\cdot \text{T-count})$ subject to the isometry constraint
verified by our \texttt{check\_isometry.py}; train on $L\in\{6,9,12\}$
pairing instances and evaluate on held-out Richardson and LMG Hamiltonians.
Baseline: the paper's construction, as reproduced here.
\end{enumerate}
